Om de helpdesk te ontzien moet er een tool komen die de webserver test. Alle domeineigenaren moeten de tool kunnen gebruiken zodat ze eerst zelf kunnen uitzoeken wat het probleem met hun site is. De te schrijven monitoringtool moet aan de gebruiker vertellen welk probleem er is. Dit doen we vanaf een ander systeem, namelijk een Windows of Linux werkstation zodat ook bij een reboot van de server de monitoringtool gestart kan worden door de gebruikers. Met deze tool kunnen gebruikers dus zelf zien waar het probleem zit en of ze zelf actie moeten ondernenem (hun domein ligt eruit), of dat ze contact moeten opnemen met de beheerders (de webserver is omgevallen), of geduld moeten hebben omdat de server aan het rebooten is.

